\documentclass[10pt,a4paper,sans]{moderncv}
\moderncvstyle{classic}
\moderncvcolor{blue}
\usepackage[utf8]{inputenc}
\usepackage[scale=0.82]{geometry}

% moderncv compose \today en anglais par défaut ; format français sans charger babel.
\renewcommand{\today}{\number\day\space\ifcase\month\or janvier\or février\or mars\or avril\or
mai\or juin\or juillet\or août\or septembre\or octobre\or novembre\or décembre\fi\space\number\year}

\name{Alexandro}{Disla}
\title{Lettre de Motivation}
\address{Rue Saint-Louis Jeanty \#14, Route de Frère}{Port-au-Prince, Haïti}{}
\phone[mobile]{(+509) 4148-3700}
\email{alexandrodisla@hotmail.com}
\extrainfo{GitHub: \url{https://github.com/AD0791}}

\begin{document}

\recipient{Ressources Humaines --- Concern Worldwide Haïti}{Pétion-Ville, Haïti\\Poste : Chargé.e Gestion de l'Information (MEAL) --- DP/CGI1210/2026}
\date{\today}
\opening{Madame, Monsieur,}
\closing{Je vous prie d'agréer, Madame, Monsieur, l'expression de mes salutations distinguées,}
\enclosure[Ci-joint]{Curriculum Vitae, attestation C.T.P.E.A}

\makelettertitle

Je vous adresse ma candidature au poste de Chargé.e Gestion de l'Information au sein du département MEAL de Concern Haïti. J'apporte cinq ans de suivi-évaluation et de gestion des données en ONG en Haïti --- la santé et le VIH à la Caris Foundation International, l'eau, l'hygiène et l'assainissement chez HANWASH, l'éducation chez Anseye Pou Ayiti --- précédés de huit ans de suivi d'indicateurs de projets publics au Ministère de la Planification. Le créole haïtien et le français sont mes langues maternelles, et je réside à Route de Frère, à proximité de Pétion-Ville.

Votre fiche décrit un système complet, du formulaire à la base des participants puis au tableau de bord : c'est celui que j'ai construit pour Anseye Pou Ayiti. J'y ai programmé les formulaires XLSForm sur ODK, avec logiques de saut et contraintes de validation pour que l'erreur soit arrêtée à la saisie plutôt que corrigée trois semaines plus tard. Les soumissions arrivent par webhook dans une base MySQL que j'ai conçue et administrée. Chaque bénéficiaire y porte une clé d'identification unique, et la base comprend une recherche de doublons, un audit des enregistrements orphelins et une sauvegarde complète. Je suis ingénieur de données en parallèle : je sais construire et maintenir ce type de système, pas seulement le spécifier.

La qualité et la protection des données sont au même niveau de priorité dans mon travail. À la Caris Foundation, j'ai établi et supervisé les protocoles de DQA en remontant les chiffres rapportés jusqu'aux registres sources, et produit les indicateurs MER destinés à l'USAID via DATIM. J'y ai aussi travaillé trois ans sur les données d'un programme VIH, où une information mal gardée expose directement une personne. J'en retiens que la confidentialité, les accès et la durée de conservation se décident dans le schéma et dans les permissions de la base, puis s'écrivent dans une procédure que l'équipe peut suivre. Pour l'analyse, je travaille sous Excel, Power BI et R, avec la formation en statistique et en échantillonnage d'un économiste appliqué.

Deux réserves, que je préfère poser franchement. Ma scolarité de quatre ans en économie appliquée au C.T.P.E.A est complétée et attestée par la pièce jointe, mais le diplôme reste en attente de la soutenance de mon mémoire de sortie. Et mon encadrement a jusqu'ici été technique : j'ai dirigé le travail des agents de collecte d'un programme multi-sites et servi de référent au personnel de suivi-évaluation, sans en être le supérieur hiérarchique. La supervision d'un.e Chargé.e Adjoint.e, évaluation des performances comprise, est la part du poste que j'apprendrai selon vos procédures.

Je reste à votre entière disposition pour un entretien.

\makeletterclosing

\end{document}
